
\subsection{定理 I 的证明}
\label{subsec:6-5-2}

本小节证明定理 I。
证明使用第 6.3.1、6.3.2 小节
给出的代数材料，
但不使用定理~\ref{thm:D} 中的谱序列。

在下一个命题中，
取以域 $R$ 为系数的诺维科夫环
$\Lambda_{0,\mathrm{nov}}
=\Lambda_{0,\mathrm{nov}}(R)$。

\begin{proposition}{6.5.1}\label{prop:6-5-1}
设 $(C,\delta)$ 为有限 d.g.c.f.z.，
并记 $C^{\{0,1\}}=C^0\oplus C^1$。
则在 $C^{\{0,1\}}$ 中存在元素
$\mathfrak z_i$（$i=1,\ldots,b$）、
$\mathfrak h_i$（$i=1,\ldots,a$）及
$\mathfrak b_i$（$i=1,\ldots,b$），
满足下列性质：
\begin{equation}
E(\mathfrak z_i)
=E(\mathfrak h_i)
=E(\mathfrak b_i)=0.
\tag{6.5.2.1}\label{eq:6-5-2-1}
\end{equation}
\begin{equation}
\begin{gathered}
\{\mathfrak z_1,\ldots,\mathfrak z_b,
\mathfrak b_1,\ldots,\mathfrak b_b,
\mathfrak h_1,\ldots,\mathfrak h_a\}\\
\text{是 }C^{\{0,1\}}\text{ 的一组标准基。}
\end{gathered}
\tag{6.5.2.2}\label{eq:6-5-2-2}
\end{equation}
\begin{equation}
\begin{gathered}
\text{存在 }\lambda_i\ge0\text{，使}
\quad\delta\mathfrak z_i
=T^{\lambda_i}e^{p_i}\mathfrak b_i,\\
p_i=\deg\mathfrak z_i\in\{0,1\}.
\end{gathered}
\tag{6.5.2.3}\label{eq:6-5-2-3}
\end{equation}
\begin{equation}
\begin{gathered}
\{\mathfrak b_1,\ldots,\mathfrak b_b,
\mathfrak h_1,\ldots,\mathfrak h_a\}\\
\text{是 }\ker\delta\cap C^{\{0,1\}}
\text{ 的一组标准基。}
\end{gathered}
\tag{6.5.2.4}\label{eq:6-5-2-4}
\end{equation}
\begin{equation}
\begin{gathered}
\{T^{\lambda_1}\mathfrak b_1,
\ldots,T^{\lambda_b}\mathfrak b_b\}\\
\text{是 }\operatorname{Im}\delta\cap C^{\{0,1\}}
\text{ 的一组标准基。}
\end{gathered}
\tag{6.5.2.5}\label{eq:6-5-2-5}
\end{equation}
\end{proposition}

\noindent\textbf{命题 6.5.1 的证明。}
令
$\mathfrak b'_1,\ldots,\mathfrak b'_b$
为 $\operatorname{Im}\delta\cap C^{\{0,1\}}$
的一组标准基，
并置 $E(\mathfrak b'_i)=\lambda_i$。
定义
\[
\mathfrak b_i=T^{-\lambda_i}\mathfrak b'_i\in C.
\]
可取 $\mathfrak z_i\in C^{\{0,1\}}$，
使
\[
\delta\mathfrak z_i
=e^{\deg\mathfrak z_i}\mathfrak b'_i
=T^{\lambda_i}e^{\deg\mathfrak z_i}\mathfrak b_i.
\]
注意 $\mathfrak b_i\in\ker\delta$，
且
$\{\sigma(\mathfrak b_i)
=\sigma(\mathfrak b'_i)\mid i=1,\ldots,b\}$
线性无关；
这里 $\sigma(\mathfrak b_i)$
是 $\mathfrak b_i$ 的首系数。
因此，由推论~\ref{cor:6-3-13}
及引理~\ref{lem:6-3-2bis}，
可取 $\mathfrak h_1,\ldots,\mathfrak h_a$，
使
$\{\mathfrak b_1,\ldots,\mathfrak b_b,
\mathfrak h_1,\ldots,\mathfrak h_a\}$
成为 $\ker\delta\cap C^{\{0,1\}}$
的一组标准基。

\begin{lemma}{6.5.3}\label{lem:6-5-3}
集合
\[
\{\mathfrak z_1,\ldots,\mathfrak z_b,
\mathfrak b_1,\ldots,\mathfrak b_b,
\mathfrak h_1,\ldots,\mathfrak h_a\}
\]
作为 $\Lambda_{0,\mathrm{nov}}$-模生成 $C$。
\end{lemma}

\begin{proof}
取 $x\in C$。
存在 $c_i\in\Lambda_{0,\mathrm{nov}}$，使
\[
\delta x=\sum_{i=1}^b c_i\mathfrak b'_i.
\]
于是
\[
\delta\left(
x-\sum_{i=1}^b
e^{-\deg\mathfrak z_i}c_i\mathfrak z_i
\right)=0.
\]
由于
$\{\mathfrak b_1,\ldots,\mathfrak b_b,
\mathfrak h_1,\ldots,\mathfrak h_a\}$
是 $\ker\delta\cap C^{\{0,1\}}$
的一组标准基，
存在 $d_i,e_i\in\Lambda_{0,\mathrm{nov}}$，
使
\[
x-\sum_{i=1}^b
e^{-\deg\mathfrak z_i}c_i\mathfrak z_i
=\sum_{i=1}^b d_i\mathfrak b_i
+\sum_{i=1}^a e_i\mathfrak h_i.
\]
这说明 $x$ 是所列
$\mathfrak z_i,\mathfrak b_i,\mathfrak h_i$
的线性组合，得证。
\end{proof}

\begin{lemma}{6.5.4}\label{lem:6-5-4}
集合
\[
\{\sigma(\mathfrak z_1),\ldots,\sigma(\mathfrak z_b),
\sigma(\mathfrak b_1),\ldots,\sigma(\mathfrak b_b),
\sigma(\mathfrak h_1),\ldots,\sigma(\mathfrak h_a)\}
\]
线性无关。
\end{lemma}

\begin{proof}
由 $\mathfrak b_i,\mathfrak h_j,\mathfrak z_k$
的选取可知，
它们的能量均为零，
且 $\sigma(\mathfrak b_i)$
与 $\sigma(\mathfrak h_j)$
组成的集合线性无关。
若结论不成立，
不失一般性，
可设 $\sigma(\mathfrak z_1)$
是
\[
\{\sigma(\mathfrak z_2),\ldots,
\sigma(\mathfrak z_b),
\sigma(\mathfrak b_1),\ldots,
\sigma(\mathfrak b_b),
\sigma(\mathfrak h_1),\ldots,
\sigma(\mathfrak h_a)\}
\]
的线性组合。
由引理~\ref{lem:6-5-3}，
后一集合便生成 $R$-向量空间
$\overline C$；这里
$\overline C\otimes_R\Lambda_{0,\mathrm{nov}}=C$。
于是
$\{\mathfrak z_2,\ldots,\mathfrak z_b,
\mathfrak b_1,\ldots,\mathfrak b_b,
\mathfrak h_1,\ldots,\mathfrak h_a\}$
作为 $\Lambda_{0,\mathrm{nov}}$-模生成 $C$。
因而
\[
\{e^{-\deg\mathfrak z_2}\delta\mathfrak z_2,
\ldots,
e^{-\deg\mathfrak z_b}\delta\mathfrak z_b\}
=\{\mathfrak b'_2,\ldots,\mathfrak b'_b\}
\]
生成 $\operatorname{Im}\delta$。
这与
$\{\mathfrak b'_1,\ldots,\mathfrak b'_b\}$
是 $\operatorname{Im}\delta\cap C^{\{0,1\}}$
的标准基这一选取相矛盾。
\end{proof}

至此，命题~\ref{prop:6-5-1} 证毕。
\hfill$\square$

\begin{remark}{6.5.5}\label{rem:6-5-5}
沿用命题~\ref{prop:6-5-1} 的记号，
有
\[
\frac{\ker\delta}{\operatorname{Im}\delta}
=\bigl(\Lambda_{0,\mathrm{nov}}^{(0)}\bigr)^{\oplus a}
\oplus
\bigoplus_{i=1}^{b}
\frac{\Lambda_{0,\mathrm{nov}}^{(0)}}
{T^{\lambda_i}\Lambda_{0,\mathrm{nov}}^{(0)}}.
\]
也就是说，$a$ 是贝蒂数，
$\lambda_i$ 是挠指数。
\end{remark}

现在回到几何情形，证明定理 I。
设 $L$ 为 $M$ 的相对 spin
拉格朗日子流形。
如第 6.4.1 小节，
利用典范模型得到有限 d.g.c.f.z.
$(C,\delta)$，
它与
\[
\bigl(C(L;\Lambda_{0,\mathrm{nov}})^+,
\delta^+_{(\boldsymbol b,b)}\bigr)
\]
上链同伦等价。
取 $R=\mathbb Q$。
可以假设
$\delta\equiv0
\pmod{\Lambda_{0,\mathrm{nov}}^+}$。

应用命题~\ref{prop:6-5-1}，
得到 $C^{\{0,1\}}$ 的一组标准基
\[
\{\mathfrak z_1,\ldots,\mathfrak z_b,
\mathfrak b_1,\ldots,\mathfrak b_b,
\mathfrak h_1,\ldots,\mathfrak h_a\}.
\]
\newpage
下面证明有关首系数
$\sigma(\mathfrak b_i)$ 的引理。

\begin{lemma}{6.5.6}\label{lem:6-5-6}
对 $i=1,\ldots,b$，有
$i_!\bigl(\sigma(\mathfrak b_i)\bigr)=0$。
\end{lemma}

\begin{proof}
令
\[
\varphi:(C,\delta)
\longrightarrow
\bigl(C(L;\Lambda_{0,\mathrm{nov}})^+,
\delta^+_{(\boldsymbol b,b)}\bigr)
\]
为保持滤过的上链映射，
它诱导同构
\[
H(C,\delta)
\cong
H\bigl(C(L;\Lambda_{0,\mathrm{nov}})^+,
\delta^+_{(\boldsymbol b,b)}\bigr),
\qquad
\overline C=H(\overline C,\overline\delta)
\cong H(L;\mathbb Q).
\]
回忆引理~\ref{lem:6-4-5} 中的映射
\[
\mathfrak p_1^{+\boldsymbol b,b}:
C(L;\Lambda_{0,\mathrm{nov}})^+
\longrightarrow
S^*(M;\Lambda_{0,\mathrm{nov}}),
\]
它满足
\[
\delta_M\circ\mathfrak p_1^{+\boldsymbol b,b}
=-\mathfrak p_1^{+\boldsymbol b,b}
\circ\delta^+_{(\boldsymbol b,b)};
\]
这里 $\delta_M$ 是
$\Lambda_{0,\mathrm{nov}}$ 系数的
通常上边界算子
\[
\delta_M:
S^*(M;\Lambda_{0,\mathrm{nov}})
\longrightarrow
S^*(M;\Lambda_{0,\mathrm{nov}}).
\]
此外，由式~\eqref{eq:6-4-9} 有
\translatornote{原书实体页 396 的式 (6.5.7)
及实体页 397 的式 (6.5.8)
把目标写成 $C(M)$；
上文的 $\mathfrak p_1$ 与 $\delta_M$
均以 $S^*(M)$ 为靶，
此处依型别统一为 $S^*(M)$，
见疑误 A-0392。}
\begin{equation}
\begin{aligned}
&\bigl(i_!-\mathfrak p_1^{+\boldsymbol b,b}\bigr)
\bigl(F^\lambda C(L;\Lambda_{0,\mathrm{nov}})^+\bigr)\\
&\qquad\subset
F^{\lambda+\lambda_0}
S^*(M;\Lambda_{0,\mathrm{nov}}).
\end{aligned}
\tag{6.5.7}\label{eq:6-5-7}
\end{equation}
这里的常数 $\lambda_0$ 定义为
\[
\lambda_0=\inf
\{\beta\cap[\omega]\mid
\beta\in\Pi(M;L),\ 
\mathcal M(\beta)\ne\varnothing,
\ \beta\ne\beta_0\}.
\]

由命题~\ref{prop:6-5-1} 证明开头
对 $\mathfrak z_i$ 的选取，
有
\[
\delta\mathfrak z_i
\equiv
T^{\lambda_i}e^{\deg\mathfrak z_i}
\sigma(\mathfrak b_i)
\pmod{T^{\lambda_i}\Lambda_{0,\mathrm{nov}}^+}.
\]
另一方面，由式~\eqref{eq:6-5-7}，
\[
\bigl(i_!-\mathfrak p_1^{+\boldsymbol b,b}\bigr)
\varphi(\delta\mathfrak z_i)
\in F^{\lambda_i+\lambda_0}
S^*(M;\Lambda_{0,\mathrm{nov}}).
\]
因此
\begin{equation}
\begin{aligned}
i_!\varphi(\delta\mathfrak z_i)
&=\bigl(i_!-\mathfrak p_1^{+\boldsymbol b,b}\bigr)
\varphi(\delta\mathfrak z_i)
+\mathfrak p_1^{+\boldsymbol b,b}
\varphi(\delta\mathfrak z_i)\\
&=\bigl(i_!-\mathfrak p_1^{+\boldsymbol b,b}\bigr)
\varphi(\delta\mathfrak z_i)
\pm\mathfrak p_1^{+\boldsymbol b,b}
\delta^+_{(\boldsymbol b,b)}\varphi(\mathfrak z_i)\\
&\equiv
\pm\delta_M\mathfrak p_1^{+\boldsymbol b,b}
\varphi(\mathfrak z_i)
\pmod{F^{\lambda_i+\lambda_0}
S^*(M;\Lambda_{0,\mathrm{nov}})}.
\end{aligned}
\tag{6.5.8}\label{eq:6-5-8}
\end{equation}
又存在 $\lambda'_0>0$，使
\begin{equation}
\delta_0\delta(\mathfrak z_i)
\equiv0
\pmod{T^{\lambda_i+\lambda'_0}C}.
\tag{6.5.9}\label{eq:6-5-9}
\end{equation}
这是因为
\[
0=\delta\delta(\mathfrak z_i)
=(\delta_0+\delta')\delta(\mathfrak z_i)
=\delta_0\delta(\mathfrak z_i)
+\delta'\delta(\mathfrak z_i),
\]
而
$\delta'\delta(\mathfrak z_i)
\in F^{\lambda_i+\lambda'_0}C$。
这里 $\delta_0$ 是
$\overline\delta:\overline C\to\overline C$
诱导的上边界算子，
$\delta'$ 是 $\delta$ 的高阶项。
由于 $\varphi$ 诱导
$H(\overline C,\overline\delta)
\cong H(L;\mathbb Q)$，
式~\eqref{eq:6-5-8}、\eqref{eq:6-5-9}
蕴含
\[
i_!\bigl(\sigma(\mathfrak b_i)\bigr)
=e^p(i_!\circ\varphi_*)
\bigl([T^{-\lambda_i}\delta(\mathfrak z_i)]\bigr)
=0\in H^p(M;\mathbb Q)
\]
（其中某个 $p\in\mathbb Z$）。
引理得证。
\end{proof}

现在可以完成定理 I 的证明。
由引理~\ref{lem:6-5-6} 得
\translatornote{原书实体页 397
把式 (6.5.10) 中的目标写成
$H^p(M;\mathbb Q)$；
但本章已定义的 Gysin 同态 $i_!$
把次数升高余维 $n$，
故此处目标应为 $H^{p+n}(M;\mathbb Q)$，
见疑误 A-0393。}
\begin{equation}
b\le
\sum_p\operatorname{rank}_{\mathbb Q}
\ker\bigl(i_!:H^p(L;\mathbb Q)
\to H^{p+n}(M;\mathbb Q)\bigr).
\tag{6.5.10}\label{eq:6-5-10}
\end{equation}
另一方面，命题~\ref{prop:6-5-1} 蕴含
\begin{equation}
a+2b=\sum_p\operatorname{rank}_{\mathbb Q}
H^p(L;\mathbb Q),
\tag{6.5.11}\label{eq:6-5-11}
\end{equation}
以及
\begin{equation}
\begin{aligned}
a=\operatorname{rank}_{\Lambda_{\mathrm{nov}}}
HF\bigl(& (L,(\boldsymbol b,b)),
(L,(\boldsymbol b,b));
\Lambda_{\mathrm{nov}}\bigr).
\end{aligned}
\tag{6.5.12}\label{eq:6-5-12}
\end{equation}
合并式~\eqref{eq:6-5-10}--\eqref{eq:6-5-12}，
便得到
\[
\begin{aligned}
&\operatorname{rank}_{\Lambda_{\mathrm{nov}}}
HF\bigl((L,(\boldsymbol b,b)),
(L,(\boldsymbol b,b));\Lambda_{\mathrm{nov}}\bigr)\\
&\quad\ge
\sum_p\operatorname{rank}_{\mathbb Q}H_p(L;\mathbb Q)
-2\operatorname{rank}_{\mathbb Q}
\ker\bigl(i_!:H^*(L;\mathbb Q)
\to H^*(M;\mathbb Q)\bigr).
\end{aligned}
\]
定理 I 得证。
\hfill$\square$
